فکر می‌کنند که ناکارآمدند، سپس این درد از بین می‌رود و ایده‌ی اینکه ما انسان باارزشی نیستیم تا حدی ناپدید می‌شود. اما در مورد افرادی که نتوانستند حس ناکارآمدی خود را کنار بگذارند چطور؟

آنها وسواس می‌کنند که با زنان بیشتری بخوابند. و مشکل اینجاست. یک استاد هنر جذب باقی مانده بود که باید او را ملاقات می‌کردم. من نمی‌خواستم از او مشاوره بگیرم برای چگونگی جذب دختران؛ می‌خواستم بدانم چگونه متوقف شوم. سپس نوعی از افرادی وجود دارند که باید در جلسات درمانی شرکت کنند.

همه در این جامعه نام او را ذکر کرده بودند. او یک حضور روحانی بود که بر دنیای جذب سایه افکنده بود، یک شخصیت اسطوره‌ای مثل اودیسه یا کاپیتان کرک یا یک HB11. او اریک وبر بود، اولین دارای مدرک حسابداری مدرن، نویسنده کتاب "چگونه باید دختران را جذب کرد" که در سال ۱۹۷۰ نوشته شد و موضوع فیلمی به همین نام بود.

او را در یک استودیوی کوچک پس از تولید ملاقات کردم جایی که فیلمی که کارگردانی کرده بود را ویرایش می‌کرد. او قطعاً جلب توجه نمی‌کرد؛ شبیه یک مدیر اجرایی تبلیغات میانسال با موی خاکستری، پیراهن آهاردار ولی دکمه‌های آن بیش از حد بالا بسته شده و شلوار سیاه بدون ویژگی بود. تنها چشمانش که با انرژی می‌درخشید نشان از آن داشت که هنوز جوان‌دلیش خاموش نشده است.

آیا از جامعه جذب آگاه هستید؟

بله، اما با حس اینکه تقلید شده‌ام به آن نگاه می‌کنم. بخش‌هایی از آنچه پس از کتاب من آمد به نظرم غیرقابل قبول بود. من به انجام کارهایی که شخص را پیچ و تاب بدهد باور ندارم. هرگز علاقه‌مند به تسخیر زنان به شکلی مستبدانه نبودم. به این علاقه‌مند بودم که کسی را پیدا کنم که او را دوست بدارم. با این حال، به جذابیت شدید علاقه‌مند نماندم. احساس می‌کردم که کارهای زیاد دیگری وجود دارد که می‌خواهم انجام دهم.

چه چیزی باعث شد بر آن غلبه کنید؟

بعد از ازدواج و کسب اعتماد به نفس بیشتری در خودم و درک این که جمع‌آوری دسته‌های زیادی در زندگی من نمی‌تواند به یأسی که دارم پایان بدهد، علاقه‌ام را از دست دادم. داشتن دو دختر که گاه به گاه مرا به داشتن نگرش جنسیتی متهم می‌کنند هم به من کمک کرد، که البته به مقدار کمی دارم.

یأسی درونی‌تان چه بود؟

فکر می‌کنم معضل وجودی این است که ما حیوانات اجتماعی هستیم، بنابراین همگی با حس ناکافی بودن کشمکش می‌کنیم. اما وقتی می‌فهمیم که آن‌قدر که فکر می‌کردیم ناکافی نیستیم و وقتی متوجه می‌شویم که همه دیگران هم همین حس را دارند، این مشکل محو می‌شود. 

آیا در آن زمان سابقه‌ای برای آن وجود داشت؟

در اواسط دهه شصت، زندگی در آمریکا به طور رادیکال تغییر می‌کرد. زنان به تازگی شروع به مصرف قرص ضد بارداری کرده بودند؛ گروه‌هایی مثل رولینگ استونز و بیتلز در صحنه حضور داشتند؛ باب دیلن بود...